\documentclass[a4paper,11pt]{article}
\usepackage[utf8]{inputenc}
\usepackage[T1]{fontenc}
\usepackage[spanish]{babel}
\usepackage{geometry}
\usepackage{array}
\usepackage{colortbl}
\usepackage{xcolor}
\usepackage{booktabs}
\usepackage{enumitem}
\usepackage{fancyhdr}
\usepackage{tikz}

\geometry{left=2cm, right=2cm, top=2.5cm, bottom=2cm}

\definecolor{violeta}{HTML}{7B2D8E}
\definecolor{azulg}{HTML}{3498DB}
\definecolor{rojog}{HTML}{E74C3C}
\definecolor{gris}{HTML}{F2F2F2}
\definecolor{verde}{HTML}{27AE60}
\definecolor{naranja}{HTML}{E67E22}

\pagestyle{fancy}
\fancyhf{}
\fancyhead[L]{\small\textbf{Nuevo Compañeros 1 — U03: La Familia}}
\fancyhead[R]{\small\textbf{Píldora formativa 3.9}}
\fancyfoot[C]{\thepage}

\setlength{\parindent}{0pt}

\begin{document}

\begin{center}
{\LARGE\textbf{PÍLDORA FORMATIVA 3.9}}\\[0.3cm]
{\Large Mediación oral: de yo a él/ella}\\[0.2cm]
{\normalsize Sección Destrezas (p.41) — Proyectable — 3 diapositivas}\\[0.2cm]
{\small Versión enriquecida secuencial para mediación MCER}
\end{center}

\vspace{0.5cm}

\noindent\fbox{\parbox{\dimexpr\linewidth-2\fboxsep-2\fboxrule}{%
\textbf{Contenido para el profesor}\\[0.2cm]
La actividad 7 (p.41) introduce por primera vez la mediación oral del MCER Companion Volume 2020. El estudiante cuenta algo especial de su familia a un compañero, y el compañero lo transmite a la clase cambiando la perspectiva (yo $\rightarrow$ él/ella).\\[0.2cm]
\textbf{Qué es mediar:} no es repetir ni traducir. Es transformar un mensaje cambiando la perspectiva para un destinatario diferente. El mediador selecciona lo esencial, transforma la persona gramatical y verifica que la información llega. Si el dato llega, la mediación es exitosa.
}}

\vspace{0.8cm}

%% ============================================================
%% DIAPOSITIVA 1
%% ============================================================

\noindent\textcolor{violeta}{\rule{\linewidth}{2.5pt}}

\section*{DIAPOSITIVA 1 — ¿QUÉ CAMBIA CUANDO CUENTAS LO DE OTRO?}

{\small Fase: Detección (awareness) — Técnica: Noticing guiado por contraste visual, pares mínimos de persona gramatical}\\
{\small\textit{Principio:} El estudiante ve dos versiones del mismo mensaje lado a lado y descubre por sí mismo qué cambia. No se da la regla antes del descubrimiento. Las palabras que cambian están resaltadas visualmente.}

\vspace{0.5cm}

\subsection*{En pantalla}

Dos columnas lado a lado:

\begin{center}
\renewcommand{\arraystretch}{1.8}
\begin{tabular}{|>{\raggedright\arraybackslash}p{6.5cm}|>{\raggedright\arraybackslash}p{6.5cm}|}
\hline
\rowcolor{azulg!15}
\textbf{LO QUE DICE MARTA:} & \textbf{LO QUE DICE EL MEDIADOR:} \\
\hline
\textbf{\textcolor{rojog}{Tengo}} un primo y una prima gemelos. & Marta \textbf{\textcolor{verde}{tiene}} un primo y una prima gemelos. \\
\hline
Se llaman Romeo y Julieta. & Se llaman Romeo y Julieta. \\
\hline
\textbf{\textcolor{rojog}{Mi}} primo toca la guitarra. & \textbf{\textcolor{verde}{Su}} primo toca la guitarra. \\
\hline
\textbf{\textcolor{rojog}{A mí me}} gusta mucho la música. & \textbf{\textcolor{verde}{A ella le}} gusta mucho la música. \\
\hline
\end{tabular}
\end{center}

\vspace{0.3cm}

Debajo de las columnas, tres preguntas:
\begin{itemize}[leftmargin=2cm, itemsep=3pt]
\item ¿Qué ha cambiado?
\item ¿Dónde está el cambio?
\item ¿Se ha perdido alguna información?
\end{itemize}

\vspace{0.5cm}

\subsection*{Instrucciones para el profesor}

\begin{enumerate}[leftmargin=1.5cm, itemsep=4pt]
\item Señale las dos columnas. Lea primero la columna izquierda completa y después la derecha. No explique nada todavía.
\item Pregunte: \textit{¿Qué observáis que cambia?} Recoja respuestas (tengo $\rightarrow$ tiene, mi $\rightarrow$ su, a mí me $\rightarrow$ a ella le). No dé la regla — deje que los estudiantes la descubran.
\item Pregunte: \textit{¿Dónde está el cambio?} (En los verbos y en los posesivos.) \textit{¿Se ha perdido alguna información?} (No. Todo llega. Pero las personas cambian.)
\item Introduzca el concepto: \textit{Esto se llama mediar. Mediar no es repetir. Es contar lo que dice otra persona cambiando la perspectiva.}
\item Haga un segundo ejemplo en vivo: diga una frase sobre usted (\textit{Tengo un hermano. Se llama Luis. Vive en Madrid.}). Pida a un voluntario que la transmita a la clase.
\end{enumerate}

\newpage

%% ============================================================
%% DIAPOSITIVA 2
%% ============================================================

\noindent\textcolor{violeta}{\rule{\linewidth}{2.5pt}}

\section*{DIAPOSITIVA 2 — LA CAJA DE HERRAMIENTAS DEL MEDIADOR}

{\small Fase: Modelado + conexión — Técnica: Tabla de transformación visual + fórmulas del mediador + plantilla de escucha activa}\\
{\small\textit{Principio:} Después del descubrimiento, el estudiante recibe las herramientas formalizadas: tabla de transformación de persona, fórmulas lingüísticas para transmitir/verificar/reparar, y plantilla de 4 casillas para la escucha activa.}

\vspace{0.5cm}

\subsection*{En pantalla}

Tabla de transformación (grande, visual, central):

\begin{center}
\renewcommand{\arraystretch}{1.5}
\begin{tabular}{|>{\centering\arraybackslash}p{5.5cm}|>{\centering\arraybackslash}p{5.5cm}|}
\hline
\rowcolor{rojog!15}
\textbf{Yo digo...} & \cellcolor{verde!15}\textbf{El mediador dice...} \\
\hline
\textbf{\textcolor{rojog}{Tengo}} & \textbf{\textcolor{verde}{Tiene}} \\
\hline
\textbf{\textcolor{rojog}{Mi}} hermano / \textbf{\textcolor{rojog}{Mi}} hermana & \textbf{\textcolor{verde}{Su}} hermano / \textbf{\textcolor{verde}{Su}} hermana \\
\hline
\textbf{\textcolor{rojog}{Me}} gusta & \textbf{\textcolor{verde}{Le}} gusta \\
\hline
\textbf{\textcolor{rojog}{Soy}} de & \textbf{\textcolor{verde}{Es}} de \\
\hline
\textbf{\textcolor{rojog}{Estudio}} & \textbf{\textcolor{verde}{Estudia}} \\
\hline
\textbf{\textcolor{rojog}{Vivo}} en & \textbf{\textcolor{verde}{Vive}} en \\
\hline
\end{tabular}
\end{center}

\vspace{0.3cm}

Fórmulas del mediador (en recuadros diferenciados):

\vspace{0.2cm}

\noindent\colorbox{azulg!10}{\parbox{\dimexpr\linewidth-2\fboxsep}{%
\textbf{PARA TRANSMITIR:} Dice que... / Tiene... / Su \_\_\_ se llama... / En su familia hay...
}}

\vspace{0.2cm}

\noindent\colorbox{verde!10}{\parbox{\dimexpr\linewidth-2\fboxsep}{%
\textbf{PARA VERIFICAR:} ¿Entendéis? / ¿Queda claro? / ¿Hasta aquí bien?
}}

\vspace{0.2cm}

\noindent\colorbox{naranja!10}{\parbox{\dimexpr\linewidth-2\fboxsep}{%
\textbf{PARA REPARAR:} No, no dice que... Dice que... / Lo que quiere decir es...
}}

\vspace{0.3cm}

Plantilla de escucha activa — MIENTRAS ESCUCHAS, ANOTA 3--5 PALABRAS CLAVE:

\begin{center}
\renewcommand{\arraystretch}{1.5}
\begin{tabular}{|>{\centering\arraybackslash}p{3cm}|>{\centering\arraybackslash}p{3cm}|>{\centering\arraybackslash}p{3cm}|>{\centering\arraybackslash}p{3cm}|}
\hline
\rowcolor{violeta!15}
\textbf{QUIÉN} & \textbf{QUÉ} & \textbf{CUÁNDO} & \textbf{DÓNDE} \\
\hline
(nombre, parentesco) & (qué hace, qué tiene) & (hora, día) & (lugar, país) \\[0.5cm]
\hline
\end{tabular}
\end{center}

\vspace{0.2cm}

\begin{center}
\fbox{\parbox{12cm}{\centering
\textbf{Recuerda:} si el dato llega, la mediación es exitosa. No importa si la gramática no es perfecta.
}}
\end{center}

\vspace{0.5cm}

\subsection*{Instrucciones para el profesor}

\begin{enumerate}[leftmargin=1.5cm, itemsep=4pt]
\item Señale la tabla de transformación. Lea cada par en voz alta: \textit{Tengo $\rightarrow$ Tiene. Mi hermano $\rightarrow$ Su hermano. Me gusta $\rightarrow$ Le gusta.} Pida que repitan cada par.
\item Señale las fórmulas. Modele con un ejemplo completo: pida a un voluntario que diga algo de su familia. Transmítalo usted usando las fórmulas, señalándolas en la pantalla.
\item Señale las 4 casillas. Explique: \textit{MIENTRAS vuestro compañero habla, anotad solo 3--5 palabras clave en estas casillas. No necesitáis escribir frases completas.}
\item Modele una segunda vez con otro voluntario. Anote visiblemente en la pizarra las palabras clave en las 4 casillas. Luego transmita a la clase usando las fórmulas y las notas. Al terminar: \textit{¿Entendéis? ¿Ha llegado la información?}
\item Señale el principio final: si el dato llega, la mediación es exitosa. No se corrige la gramática durante la mediación.
\end{enumerate}

\newpage

%% ============================================================
%% DIAPOSITIVA 3
%% ============================================================

\noindent\textcolor{violeta}{\rule{\linewidth}{2.5pt}}

\section*{DIAPOSITIVA 3 — PRACTICA: CUENTA ALGO DE TU FAMILIA}

{\small Fase: Aplicación anticipada + verificación + autorregulación — Técnica: Mediación oral con banco de historias + weaning off + autoevaluación}\\
{\small\textit{Principio:} El estudiante practica la mediación con andamiaje máximo (tabla + fórmulas + banco de historias) y luego con andamiaje reducido (weaning off). El banco de 8 historias curiosas resuelve el bloqueo de contenido.}

\vspace{0.5cm}

\subsection*{En pantalla}

\begin{center}
{\Large\textbf{CUÉNTALE A TU COMPAÑERO UNA COSA ESPECIAL DE TU FAMILIA.}}\\[0.2cm]
{\large Tu compañero se lo va a contar a la clase.}
\end{center}

\vspace{0.3cm}

Ejemplo del libro: \textit{Tengo un primo y una prima gemelos. Se llaman Romeo y Julieta...}

\vspace{0.3cm}

\noindent\colorbox{gris}{\parbox{\dimexpr\linewidth-2\fboxsep}{%
\textbf{BANCO DE HISTORIAS PARA INSPIRARTE} (si no se te ocurre nada, adapta una con tus datos):\\[0.2cm]
\begin{enumerate}[leftmargin=1cm, itemsep=2pt]
\item Mi tía tiene dos hijos gemelos. Se llaman Blanco y Negro. Blanco estudia Música y Negro estudia Matemáticas. Tienen 14 años.
\item Mi abuela tiene 68 años y juega al fútbol todos los sábados a las 10 de la mañana. Su equipo se llama Las Abuelas. Ella es la número 7.
\item Mi primo se llama Leo. Su padre es argentino y su madre es japonesa, pero viven en España. Leo habla tres idiomas en casa.
\item En mi familia todos tocan un instrumento. Mi padre toca la guitarra, mi madre toca el piano, mi hermano toca la batería y yo canto.
\item Mi hermano mayor tiene 17 años y estudia todos los días a las 5 de la mañana. Luego desayuna, corre 3 kilómetros y llega al instituto a las 8.
\item Tengo 12 primos. Viven en 5 países diferentes: México, Francia, Marruecos, Italia y Alemania.
\item Mi hermana tiene 40 plantas en su habitación. Tienen nombres de personas: Rosa, Margarita, Jazmín...
\item Mi abuelo cocina todos los viernes para toda la familia. Cocina para 15 personas. Su especialidad se llama arroz con pollo.
\end{enumerate}
}}

\vspace{0.3cm}

Esquema de la actividad en 3 pasos:

\begin{center}
\renewcommand{\arraystretch}{1.5}
\begin{tabular}{|>{\centering\arraybackslash}p{1cm}|>{\raggedright\arraybackslash}p{11cm}|}
\hline
\rowcolor{verde!15}
\textbf{1} & A cuenta a B (2 min) — B escucha y anota en las 4 casillas. \\
\hline
\textbf{2} & B cuenta a A (2 min) — A escucha y anota. \\
\hline
\rowcolor{violeta!10}
\textbf{3} & Puesta en común: cada mediador presenta lo de su compañero a la clase con las fórmulas. \\
\hline
\end{tabular}
\end{center}

\vspace{0.3cm}

\noindent\textbf{Weaning off:}
\begin{itemize}[leftmargin=2cm, itemsep=2pt]
\item Primera ronda: tabla de transformación y fórmulas visibles en pantalla.
\item Segunda ronda (si hay tiempo): tabla oculta, solo fórmulas visibles.
\end{itemize}

\vspace{0.3cm}

\noindent\textbf{Autoevaluación} (30 segundos, al final):
\begin{enumerate}[leftmargin=2cm, itemsep=2pt]
\item ¿Mi compañero ha entendido lo esencial?
\item ¿He usado las fórmulas de la tarjeta?
\item ¿Ha sobrado o faltado información?
\end{enumerate}

\vspace{0.5cm}

\subsection*{Instrucciones para el profesor}

\begin{enumerate}[leftmargin=1.5cm, itemsep=4pt]
\item Reparta las tarjetas de estrategia de mediación (Caja 3, tarjeta D) si no lo ha hecho ya.
\item Señale el banco de historias. Diga: \textit{Si no se os ocurre nada, podéis adaptar una de estas historias con vuestros datos. También podéis inventar.}
\item Forme parejas. Indique: A cuenta algo especial de su familia a B. B escucha y apunta solo palabras clave en las 4 casillas. Dé 2 minutos. Luego intercambian.
\item Puesta en común: cada mediador presenta lo de su compañero a la clase. Señale las fórmulas: \textit{Dice que tiene... / Su \_\_\_ se llama... / En su familia hay...}
\item Al terminar cada mediación, pregunte a la clase: \textit{¿Ha llegado la información? ¿Se ha perdido algo?}
\item No corrija la gramática durante la mediación. Después, use reformulaciones correctivas: si un estudiante dice \textit{su hermano se llaman Carlos}, repita \textit{Ah, su hermano se llama Carlos. Muy bien.}
\item Si hay tiempo, proponga una segunda ronda sin la tabla visible (weaning off).
\item Cierre con la autoevaluación: cada estudiante responde las 3 preguntas en 30 segundos.
\end{enumerate}

\end{document}
